% !PRINTED 400
\weberpage{400}
%<<CONT
the positive values of $t$, $u$; a second is positive and
smaller than 1, and the two others are negative. If $T$, $U$
is the smallest positive solution of the Pell equation, then
\[
\Theta = \frac{T + U\sqrt{D}}{2}
\]
is the smallest improperly fractional positive unit belonging to $\sqrt{D}$,
and we can show that all units belonging to $\sqrt{D}$ are contained
in the form $\pm \Theta^n$, in which $n$ denotes an integer. For this
it suffices to show that the expression $\Theta^n$, for a positive $n$,
yields all the units that are greater than 1.
But this is very simple; for if $\Theta_1$ is any one of these
units, then it will lie between two successive
powers of $\Theta$, since these powers grow to infinity with the
exponent; thus
\[
\Theta^n \leqq \Theta_1 < \Theta^{n+1},
\]
hence
\[
1 \leqq \Theta_1 \Theta^{-n} < \Theta.
\]
But since $\Theta_1 \Theta^{-n}$ is likewise a unit, it follows, unless
the equality sign holds, that $\Theta_1 \Theta^{-n}$ is greater than 1 and smaller than
$\Theta$, which is contrary to the assumption that $\Theta$ is the smallest
improperly fractional unit; hence we must have
\[
\Theta_1 = \Theta^n
\]
which was to be proved.

\gesperrt{Remark concerning the \gesperrt{Gauss} theory of quadratic
forms}\footnote{Cf. Disq. ar. art. 183 ff., \gesperrt{Dirichlet}-\gesperrt{Dedekind}, Vorlesungen
über Zahlentheorie. Fourth edition. \S.\,72 ff.}.

In order to establish the connection of our considerations with
the otherwise well-known \gesperrt{Gauss} theory of quadratic
forms, we add a few remarks, which in our
context are not strictly required, and which may therefore
also be passed over.

Two quadratic forms,
\[
(A, B, C), \quad (-A, B, -C)
\]
in the notation of \gesperrt{Gauss}, are equivalent to one another in
the sense of \gesperrt{Gauss} then and only then, when the period
%CONT>>
